% ============================================================
%  R. Nielsen, "On S. Kierkegaard's ``Mental State''"
%  [Om S. Kierkegaards „mentale Tilstand“]
%  Nordisk Universitets-Tidsskrift, 4th year, 3rd number
%  (Copenhagen 1858), pp. 1--29.
%
%  TRANSLATION (English)
%  Translated from transcription.tex (Danish) in this directory.  Method:
%  ../../../TRANSLATION-PLAYBOOK.md.  No published translation of this
%  article is known.  Nielsen quotes Kierkegaard at length (the Christian
%  Discourses, For Self-Examination, Practice in Christianity, This Must Be
%  Said, The Moment, the Fædrelandet articles, the Postscript); every
%  quotation is translated from NIELSEN'S Danish as printed, which is the
%  copy-text here.  No published English translation of Kierkegaard was
%  consulted or substituted, and no quotation was corrected toward
%  Kierkegaard's own text; where Nielsen's print differs from it, the
%  divergence is recorded in a % TR: comment.
%
%  ---- REGISTER ----
%  Scholarly, moderately literal, readable English; American spelling.
%  Nielsen's long periods are broken only where English will not carry
%  them, never across a page marker or a footnote anchor.
%
%  ---- TERMINOLOGY ----
%  "mentale Tilstand" = ``mental state'', always inside Nielsen's
%  quotation marks (the phrase is his opponent's, from the epigraph);
%  "sikker Dom" = sure judgment.  Sandhedsvidne = witness to the truth;
%  Christendom = Christianity, Christenheden = Christendom; det
%  Bestaaende = the established order; Forfatter-Virksomhed = authorship;
%  Inderliggjørelse = inwardization; opbyggelig = upbuilding; Forargelse
%  = offense; Existensmeddelelse = existence-communication; Videnskab =
%  science; Tro = faith.  Kierkegaard's works in their standard English
%  titles; all other titles in their original language.
%
%  ---- CONVENTIONS ----
%  - Page markers \opage{N} and "% --- p. N ---" at the same joints as in
%    transcription.tex; a word the Danish divides across a page turn is
%    kept whole in English.
%  - The same paragraphs, \emph spans (the article has no italic: every
%    \emph is letterspacing in the print), the two *) footnotes (pp. 17,
%    25), the epigraph, rules and centre blocks as the Danish.
%  - Quotation marks „...“ -> ``...'', inner quotations `...'.  The print
%    reuses „...“ for nested quotations and often leaves a mark unclosed;
%    the English closes every quotation and records each site in % TR:.
%  - Ellipses keep the printed number of spaced points.
%  - German stays German, as printed.
%
%  ---- DEFECTS: DANISH AS PRINTED, ENGLISH CORRECT ----
%  The transcription carries every printer's error; the translation
%  renders the evident sense and logs it at the site: p. 1 words dropped
%  ("…ogsaa Magt"), supplied in brackets; p. 2 a clause without its verb;
%  p. 10 "charakterisk"; p. 13 "i" doubled; p. 22 "christendum" (= christen
%  dum); p. 28 "störe Gud"; p. 29 "Forudsætninger.,", "komme" (= komne).
%
%  ---- HOW THE TRANSLATION WAS CHECKED (2026-09-23) ----
%  (1) Three translators (pp. 1--10, 11--19, 20--29), each re-reading its
%  English against the Danish sentence by sentence.  (2) Mechanical audit
%  (tr_audit.py): per printed page the same \opage joints, \emph spans,
%  footnotes, rules, centre blocks and paragraph breaks as
%  transcription.tex, and an English/Danish word ratio within 0.95--1.60:
%  29 pages and the head, 0 flagged.  (3) Independent review: three readers
%  who had not translated the pages read every sentence against the
%  Danish; 17 corrections applied, among them p. 2 a conditional read as
%  two questions; p. 17 "paa det Sidste" = of late; p. 19 "Følgeblade" is
%  Kierkegaard's own heading (Supplementary Leaves); p. 20 "orienteret";
%  p. 23 "altsaa" dropped twice from the Apostle's replies.  (4) Six
%  readings the translators doubted were checked at the image (600 ppi):
%  four confirmed as printed and now logged in the transcription, and two
%  paragraph breaks at page turns (before pp. 25 and 27) found missing
%  from the transcription and set in both files.
%
% ============================================================
\documentclass[12pt,a4paper]{article}
\usepackage[utf8]{inputenc}
\usepackage[T1]{fontenc}
\usepackage[english]{babel}
\usepackage[protrusion=true,expansion=true]{microtype}
\usepackage[margin=1.3in]{geometry}
\usepackage{setspace}
\usepackage{libertinus}
\usepackage{libertinust1math}
\usepackage{fancyhdr}
\usepackage[colorlinks=true,linkcolor=black,urlcolor=black,
  pdftitle={On S. Kierkegaard's "Mental State"},
  pdfauthor={R. Nielsen}]{hyperref}

\setlength{\headheight}{15pt}
\pagestyle{fancy}
\fancyhf{}
\fancyhead[L]{\textit{On S.\ Kierkegaard's ``Mental State''}}
\fancyhead[R]{\thepage}
\setlength{\parindent}{1.5em}
\setstretch{1.3}

% the print marks both its footnotes *) and restarts per printed page
\renewcommand{\thefootnote}{*)}

% ---- original-page markers ----
\usepackage{marginnote}
\newcommand{\opage}[1]{\marginnote{\footnotesize\textit{[#1]}}}
\begin{document}

% --- p. 1 ---
\opage{1}
\begin{center}
  {\large On S.\ Kierkegaard's ``Mental State.''}

  \medskip
  By Prof.\ R.\ Nielsen.
\end{center}

\bigskip
\begin{center}\rule{0.25\linewidth}{0.4pt}\end{center}
\bigskip

\begin{quote}\small
% TR: epigraph quoted by Nielsen from Nyt Theologisk Tidsskrift 8:1 (1857); the journal title stays in Danish.
. . . . . Among the strangest things we have witnessed in our
day was not so much this, that a man about whose mental
state in the last period of his life one cannot have any
sure judgment sought a Herostratic fame by hurling insults
and accusations against a man's character . . . ., but
even more wondrous was. . . . .\quad Nyt Theologisk Tidsskrift, edited
by Dr.\ C.\ E.\ \emph{Scharling} and Dr.\ C.\ T.\ \emph{Engelstoft}, 8th volume,
1st number, 1857, pp.\ 203 and 204.
\end{quote}

\bigskip
\begin{center}\rule{0.25\linewidth}{0.4pt}\end{center}
\bigskip

% TR: "mentale Tilstand" = ``mental state'', kept in Nielsen's quotation marks throughout (the phrase is his opponent's, from the epigraph).
In order to be able to have ``any sure judgment'' about S.\ Kierkegaard's
``mental state,'' whether in the first or the
last period of his life, it is an indispensable condition
that one should at least with some reflection consider
what it is, then, taken as a whole, that ``this man''
really wanted, what powers he unfolded, what means
he employed, what weapons he used.

% TR: p. 1 "Anstrengelsen" (with e, beside "Aandsanstrængelser") is so printed (checked at the image, 600 ppi); same sense.
% TR: printer's error (noted by the transcriber): the print reads "er jeg saa rigtignok ogsaa Magt", words having dropped out (Kierkegaard: "ogsaa i Besiddelse af en Magt"). Rendered by the evident sense, the missing words supplied in brackets.
``In torments such as a human being has seldom experienced,
in exertions of spirit that in the course of 8 days would probably
rob another of his reason, I am then, to be sure, also [in possession of a]
power, undeniably a consciousness seductive for a poor human being,
if the torment and the exertion were not to such a
degree the preponderant thing that often enough my wish is
% --- p. 2 ---
\opage{2}
death, my longing the grave, my desire that my wish
and my longing might soon be fulfilled.''

He who thus himself announces his coming forward must,
if he is otherwise in his right mind, surely have a burden
on his conscience; must, if he really knows
what he is saying, have something to say in earnest. What,
then, does he want? ``Honesty!''

And what is it, then, that he has to say?

% TR: p. 2 "hvilket Dit Liv forresten, min Ven" has no verb in the print (checked at the image; Kierkegaard has "forresten er"); rendered with the verb understood.
\emph{``This must be said, so let it then be said:
Whoever you are, whatever your life may otherwise be, my
friend---by ceasing (if you do otherwise take part in it) to
take part in the public worship of God
as it now is (with the claim of
being the Christianity of the New Testament) you have
continually one, and a great, guilt the less: you
do not take part in making a fool of God by
calling the Christianity of the New Testament
what is not the Christianity of the New Testament.''}

These are strong words, a judgment that awakens horror. If
he who dares to speak thus, and to judge thus, now also
really one who has the right, is not society then threatened with a
crisis of spirit that in its consequences could become worrying for
the state, dangerous for the spiritual peace of the people?

``Herewith---yes, O God, let now be done what Your will is,
infinite love!---herewith I have spoken.''

This is, to be sure, religious; but is it not, perhaps,
once again all too religious?

``Perhaps one person is to hear it in such a way that he does what
I say, another in such a way that he takes it as
well-pleasing to God, thinks to render God a service of worship by
joining in raising an outcry against me: neither of these
concerns me; what concerns me is only that it shall be said.''

In all this there is a resignation, a resoluteness,
a pathos, which undeniably points to a quite peculiar
``mental state''; but which? One listens, and
% --- p. 3 ---
\opage{3}
% TR: p. 3 "Öjeblik" (j for i) is so printed (checked at the image) = Öieblik, moment.
asks with anxiety: which? If now, at so
critical a moment, the most clear-sighted of the age succeed in seeing through
the ``mental state'' of the man who has come forward, and in convincing themselves
that it is by no means one who has the right that they have before them,
but a feeble-minded Herostratus, a poor genius who in
his overexcitement has gone so far that he has gone out of
his mind; then the contemporary age grasps at once that the man is mad, and
it is a great comfort to the age ``that the man is mad.'' For, it
goes without saying, when the man is mad, then the tension is
ended, the crisis over, and ``honesty'' falls away of
itself.

The clear-sighted, who speculate in the advantage of spreading
and maintaining the suspicion against K.'s ``mental state,''
so beneficial to the peace of souls, are then naturally
also clear with themselves about how
this suspicion can find its confirmation from more than one
point of view.

From the one side one can concede the truth of
everything K.\ says about the difference between the Christianity
of the New Testament and the Christianity of the present day; one
can admire the acumen, clarity, consistency
and certainty with which he knows how to make the dissimilarity
evident, and still find his ``mental state'' lamentable.
His derangement of mind then consists by no means
in this, that he does not see what Christianity originally is,
but in this, that in dead earnest he wants to force upon us ``the
Christianity of the New Testament,'' wants to screw the age
back to that historically superseded standpoint; ``so,
still that first unconditional faith in the God-man's
miracles, still that first ascetic hatred of the world, that
dark view of life: the man is surely mad!''

From the other side, on the contrary, one begins with
the concession that the Christianity of the New Testament
is the only Christianity valid for all times.
K.'s error, then, is not that even in our day he
demands that the Gospel be acknowledged; but the delusion is that
% --- p. 4 ---
\opage{4}
he does not rightly know the Gospel, since he does not see
that it is precisely the Christianity of the New Testament that
at present finds its most essential expression, its best conditions
and most necessary guarantees in the established order.
``Unhappy feeble-mindedness, that he could not after all see
that it is only something external that has changed here, not the
internal; not see that congregational life is still marked
by the same brotherly love as in the New
Testament; was not granted the grace to see that the servants of the Word,
if it came to that, are at once willing to forsake
all for Christ's sake, and thus still, save for a few
small changes, resemble the disciples in the New Testament;
the poor blinded man!''

The circumstance that one can thus by two different
ways arrive at the same result concerning K.'s
``mental state'' one must not, however, be hasty
to interpret as if we had here two proofs that fused
together into mutual reinforcement. The standpoints indicated
are indeed so contradictory to each other that what
is adduced from the one side as proof of veracity
must be adduced from the other side as a sign of insanity,
an inner contradiction that is not exactly fitted to
make the judgment secure.

No, in order to have a sure judgment about K.'s ``mental
state'' in the last period of his life, it is necessary
that one first form a sure judgment about his
``authorship'' in the first period. This judgment
has in self-complacent superficiality hitherto kept itself so
vague, and herein lies the reason why the misunderstanding
at the last had to grow to such a height.

To be sure, K.\ aroused attention at his first appearance,
but what was so highly peculiar in his style and manner of writing
at once brought it about that he had to repel some
to the same degree as he attracted others. At
another time, more favorable to deeper-going criticism,
this might perhaps have had the beneficial
% --- p. 5 ---
\opage{5}
consequence that his view had at once been so strongly lit
up from the opposite standpoints that the general
consciousness could have followed from the start; now,
on the other hand, the interpretation was on the whole left to individual
discretion, criticism went about in negligee: the various
opinions held good according to taste and fancy.

Thus, to keep here to the religious, some
thought that K.\ really lacked earnestness, while others
found in him the strictest earnestness. Among these last
stood even Bishop Mynster himself. Already the first
``upbuilding discourses'' made so earnest an impression on
the old Bishop that he, while pronouncing his recognition with weight,
even feels moved to give a
psychological contribution to a closer understanding of the young
man's strange earnestness.

% TR: the print does not reopen the quotation after the parenthesis ("noget Rörende" has no opening mark); the English supplies it, so the pair balances.
``There is for me,'' says Bishop Mynster (Kirkelig Polemik.
Intelligensblade, ed.\ J.\ L.\ Heiberg. Jan.\ 1844.)
``something touching in Magister Kierkegaard's always dedicating his
upbuilding discourses to the memory of his late father. For I
too knew this worthy man; he was a plain and simple
citizen, went unpretentiously his quiet way through
life, had never gone through any
`philosophical bath' whatever. How does it happen, then, that the son,
in his rich cultivation, as often as he wishes to write down upbuilding
discourses, always turns his thought to that man who
long ago entered into rest? Whoever has read
the lovely discourse---or let us just call it a sermon---`The
Lord gave, the Lord took, blessed be
the name of the Lord,' will understand it. The son has, like me, seen
the old father at bitter losses, he has seen him fold
his hands, bow his venerable head, he has heard
his lips utter that word, but at the same time seen his whole
being utter it in such a way that he felt
what he so beautifully develops about Job, that `he too is
a teacher of mankind who had no teaching to hand
over to others, but left the race himself as a
% --- p. 6 ---
\opage{6}
pattern, his life as guidance for every human being'
who saw it; he has felt that the old man had in his
pious word overcome the world. And what the son there
in the house of mourning saw in his old father, that he wrote down
in a sermon that will speak to, refresh every feeling
heart.'' . . . .

If the general judgment about K.'s ``religious earnestness''
was already somewhat uncertain, the scientific judgment
about his religious development of concepts became even entirely confused;
he called faith a paradox, after all, and made use
of ``the indirect form of communication'': he repelled the
speculative thinkers, he repelled the theologians.

That the rejection of the indirect form was nevertheless still
as little unanimous as offense at ``the paradox''
was general can be seen expressly, among other things, from
the way in which Bishop Mynster in the above-mentioned article
also speaks of ``Fear and Trembling.'' ``I have also,''
says the Bishop, ``read the remarkable writing `Fear and
Trembling,' and whatever I may have missed in it,
certainly not a deep religious ground, not a mind that
is equal to facing life's highest problems;
% TR: "træde ... under Öine" = to face (lit. step before the eyes of); used consistently pp. 6--8.
% TR: the print never closes the Jacobi quotation; the English closes it after "u. s. w.", before the reference. "S. 32" rendered "p. 32"; the German stays as printed.
it has vividly reminded me of the famous passage
in Jacobi: `Ja, ich bin der Atheist und Gottlose, der
dem Willen, der Nichts will, zuwider, lügen will, wie Desdemona
sterbend log, lügen und betrügen will, wie der
für Orest sich darstellende Pylades u.\ s.\ w.' (Jacobi an
Fichte p.\ 32), of which it is, however, by no means an imitation
or echo.'' By this turn of phrase, then,
the Bishop himself lays weight on the indirect communication,
and rejects those who would accuse the author of ``untimely
hunting after the piquant,'' when, after having
expressed himself on the book's title, he expressly adds:
``Thus this struggle in fear and trembling has not
the least in common with the complacent coquetry with
which the moral geniuses strive to get into such
situations where they can show that they have not in vain
% --- p. 7 ---
\opage{7}
% TR: p. 7 the German is printed "zu handlen" (checked at the image); kept as printed.
% TR: the print closes only the inner quotation after "handlen"; Mynster's quotation (opened p. 6 at "Saaledes") is never closed. The English closes both.
read Schelling, that they too are in a position `einzig göttlich
zu handlen.'\,''

Had Bishop Mynster now really been able to retain
this first impression of his, that here in ``Fear and Trembling''
he met ``a deep religious ground,''
% TR: the print has a closing mark after "Öine," but no opening mark for "et Sind ..." (Mynster's words, p. 6); the English supplies the opening mark.
``a mind that was equal to facing life's highest problems,''
had he in this respect had confidence in his
own judgment, had himself faced these ``highest problems,''
and so by his example become a guide
for the many who so gladly leaned on
his authority: what an entirely different turn would the course of the matter
then have taken from the very beginning!
When one grants an author his presuppositions,
then it is surely self-evident that one must also grant
him the necessary consequences of the presuppositions, that is to say,
if the concession rests on insight and not on a
fleeting mood. But in ``Fear and Trembling'' we have, after all,
the presuppositions of principle. It belongs precisely to
what is peculiar in Kierkegaard's nature that from
the beginning he is so at one with himself, so inwardly
finished. He was not a young man who grew old with the years,
not a merry one who later became so earnest,
not an aesthete who later became religious; no,
he was originally all that he was, in a peculiar
doubling, in his youth old, in his jest earnest,
in his pain glad, in his severity mild, in his
bitterness wistful: Kierkegaard is to such a degree an a priori
nature that he almost lacks the perfectible.

In ``Fear and Trembling'' the principle is stated, completely
stated; we see in this writing as in a mirror the author's whole intellectual
physiognomy. If there is now in this physiognomy
something essentially distorted, then it is truly striking
that a man like Bishop Mynster could, through these distorted
features, have thought he saw---``a deep religious
ground.'' Or, if there is here in this disposition something false,
and this falseness, as indeed it must, insofar as it
% --- p. 8 ---
\opage{8}
lies in the presuppositions and proceeds from the principles,
is even a fundamental falseness; is it not then inexplicable that
a connoisseur like Bishop Mynster, instead of at once catching
sight of this fundamental falseness, could let himself be dazzled,
even to the degree that behind the pseudonymous character-mask
he believed he discovered a spirit akin to
his own, ``a mind that was equal to facing life's
highest problems''?

Restlessly and with overflowing productivity K.\ worked
for a series of years at clarifying what he called
his ``existence-problems.'' His indirect communication,
which behind its veil preserved inwardness, softened the
polemical, and set a reconciling separation between actuality
and ideality, procured quiet, so that the tactical
advance could take place in all stillness.

It was Kierkegaard's plan to carry through a revision
of the categories of existence, and to separate them from the
muddle into which a misguided scientificality had
dragged down ``life's highest problems.'' He has solved
this task with a certainty, thoroughness and many-sidedness
that will in times to come warrant him a place
among the thinkers of the first rank. But despite his respect
for true scientificality where science applies,
and although in practice he is a systematician second to none,
it nevertheless belonged essentially to the sharpening of the problems of inwardness
that he had set himself, continually, even
systematically, to break the systematic form, everywhere to repel
the stupidly scientific. He employed this tactic
with great dexterity, added building to building, and enjoyed
at last the satisfaction that the topping-out wreath was already waving
on the last building of the complex when paragraph-earnestness came
along, and naturally regarded the whole as a mass of ``stray
thoughts, aphorisms, fancies and flashes.''

That Kierkegaard wanted to do away with ``speculation,'' and that
it was ``the system'' above all that suffered, was so conspicuous
a fact that there were even those who thought that
% --- p. 9 ---
\opage{9}
he passed off faith as a paradox merely to tease ``the system.''
Of this one can only say that everyone judges
as he has understanding to, and besides recall the motto of the ``Stages'':
``Solche Werke sind Spiegel, wenn ein Affe
hinein guckt, kann kein Apostel heraus sehen.''

How important it must have been to K., purely
practically and for the sake of existence, to get ``speculation''
separated from ``faith'' will be understood when one
looks at the intimate connection of the dogmatic-speculative school
with what modern philosophy since Spinoza
has called ``immanence.'' In immanence
not only is pagan ethics, in its noblest, Socratic
form, grounded, but even within Judaism and Christianity
there is, according to K.'s understanding, a popular, humane
appropriation of the revealed religion that in a
certain way likewise corresponds to immanence. So far is
K.\ from wanting this religious-ethical intermediate
stage overlooked, rejected or disdained that on the contrary
he has striven to develop its significance in a series
of ``upbuilding discourses.''

For ``a sure judgment'' about K.'s authorship
it is necessary that emphasis be laid here once more on
this side of his productivity, which during the recent disputes has been all too much
overlooked; for these discourses have
just that mildness, soberness, quietness, that awakening and
yet calming quality which one otherwise thinks is missing in
K., and in which one is then accustomed to place the genuinely evangelical.
What a difference it makes in this respect whether
in one's judgment one is led by immediate feeling or by a
clear thought, a determining category, can be seen
from the following. Among the upbuilding discourses belongs precisely
also the one on Job's words: ``The Lord gave, the Lord took,
blessed be the name of the Lord;'' this discourse Bishop
Mynster, as already mentioned, found so upbuilding that it
must be able to ``speak to, refresh every feeling heart,''
and would on that account call it a sermon. Here
% --- p. 10 ---
\opage{10}
% TR: printer's error (noted by the transcriber): "charakterisk" for "charakteristisk"; rendered "characteristic".
is a characteristic trait. The famous preacher appeals
to ``every feeling heart,'' and wants the
upbuilding discourse on Job called a ``sermon,'' whereby
he indicates that the view of life developed in the discourse
is on a level with that in the ecclesiastical address;
the dialectical author, whose endeavor it is to get
the various categories of existence clarified, on the other hand
does not want his discourse called a sermon, because he
wants ``sermon'' to correspond to the Christian, and knows
within himself that this discourse of his does not express what is in
the proper sense Christian. His rejoinder through
% TR: "Uvidenskabelig Efterskrift" (Nielsen's short title) = ``Unscientific Postscript''; "S. 204 og 205" = pp. 204 and 205.
the pseudonym in ``Unscientific Postscript'' pp.\ 204 and
205 shows this. ``Magister Kierkegaard has presumably known well enough
what he was doing when he called the upbuilding
discourses: upbuilding discourses, and why he refrained
from the use of Christian-dogmatic determinations, from
naming the name of Christ, etc., which one otherwise in our time
does quite freely, while the categories, the thoughts, the dialectical
in the presentation are only those of immanence. As the
pseudonymous writings, besides what they directly are, indirectly
are a polemic against speculation, so too are
these discourses, not by not being speculative,
but by not being sermons.''

One will easily understand that there is no question here of ``a
Christian severity over against a given Christian mildness'';
what is at issue here is to find ``honesty'' in the principles,
to vindicate for the Gospel its distinctiveness, and to prevent
its being categorially confused with a related
but nevertheless subordinate standpoint of resignation and piety.
Those who have followed the Kierkegaardian movement with some attention and reflection
must surely know that this author did not let the matter rest at
having some pseudonyms muster a brilliant dialectic
in order to exhibit ``Stages on Life's Way,'' and with formal
acumen maintain the claim that faith is a paradox;
% TR: SLICE ENDS mid-sentence (continued at the head of p. 11). Full Danish sentence across pp. 10/11: "De, som ... maae jo vide, at denne Forfatter ikke lod det beroe ved at lade nogle Pseudonymer opbyde en glimrende Dialektik ... og med formel Skarpsindighed hævde den Paastand, at Troen er et Paradox; men, at det i den Grad var ham om Sagens rette | Forstaaelse og religiös-praktiske Gjennemförelse at gjöre, at han i en rig Mangfoldighed af populære Foredrag: [titles] o. s. v., har paa det Omhyggeligste viist, og ... belyst, hvad Kategorierne gjælde, ..." This part renders p. 10's words "men, at det i den Grad var ham om Sagens rette" as "but that he was concerned to such a degree with the right", deferring "Sagens" (of the matter) to the next part, which should continue: "understanding and religious-practical carrying-through of the matter that he, in a rich multiplicity of popular addresses: ..., has most carefully shown ...".
but that he was concerned to such a degree with the right
% --- p. 11 ---
% TR: The slice opens mid-sentence (p. 10 ends "men, at det i den Grad var ham om Sagens rette"; p. 11 continues "Forstaaelse og religiös-praktiske Gjennemförelse at gjöre, at han ..."). Rendering of this part: "understanding and religious-practical carrying-through of the matter that, in a rich variety of popular discourses ..., he has shown ...", to follow "... but that he was concerned to such a degree with the right".
% TR: "hvad Kategorierne gjælde": "what the categories amount to", i.e. what is at stake in them.
\opage{11}understanding and religious-practical carrying-through of the matter
that, in a rich variety of popular discourses:
\emph{Upbuilding Discourses in Various Spirits, Works of Love,
Christian Discourses, Discourses at the Communion on Fridays}, etc.,
he has shown with the utmost care, and illuminated in the most varied
shadings, what the categories amount to, how, if possible, the people
were to be worked upon, how one ought to speak in the congregation.

% TR: "forsaae sig paa ... oversaae": the Danish plays on "forsaae sig" (become infatuated) and "oversaae" (overlooked); not reproducible.
``The congregation! But is it then really not true, what is nevertheless
asserted from so many sides, that S.\ Kierkegaard only wanted to be
`that single individual,' and became so infatuated with `the single
individual' that he entirely overlooked the congregation?'' No, it is
really untrue, categorically untrue. Only by a heightened emphasis on
`the single individual' did it become possible for him to secure the
concept of the congregation against confusion with communities of a
lower order; ``insofar,'' he says, ``as there is, religiously, a
`congregation,' this is a concept that lies on the other side of `the
single individual,' and that above all must not be confused with what
may have validity politically: the public, the crowd, the numerical,
etc.''

% TR: "at ogsaa . . . De blive flere": the ellipsis breaks the construction as printed; "De ... som" read as "those who", not the polite "you".
``To me, the upbuilding author, it was and is therefore a joy that
also . . . those grow more numerous who become attentive to what is said
about `the single individual'; it was and is a joy to me, for
certainly I have faith in the rightness of my thought in spite of the
whole world, but what, next to that, I would give up last is my faith
in individual human beings. And this is my faith: that as much
confusion and evil and abomination as there can be in human beings as
soon as they become the irresponsible and unrepentant `public,'
`crowd,' and the like, just as much truth and goodness and lovableness
is there in them when one can get them as single individuals. Oh, and
to what degree would human beings not become --- human beings, and
lovable, if they would become single individuals before God!''
(``On My Work as an Author,'' p.~12).
% --- p. 12 ---
\opage{12}The exclusive interest with which K.\ worked for the religious
showed itself, then, also in this, that he not only wrote sermons but
also, now and then, when occasion offered, actually preached,
preferably at the Friday communion. Why he preferred Friday in
particular he once explained in conversation in his joking way,
lamenting that he did not have the voice to be a Sunday pastor and be
heard by so many, yet thinking that with his weaker voice he might well
be suited to be a Friday pastor.

For a closer determination of the point of view under which K.\ wished
the Christian discourse, in the proper sense, to be understood, we shall
here refer to the third part of \emph{Christian Discourses}, and bring
out some points from its first discourse: ``\emph{Guard your foot when
you go to the house of the Lord.}'' The discourse begins thus:

% TR: The print closes this quotation on an ellipsis with no closing quotation mark; the English supplies it.
``How still, how safe everything is in the house of God. For one who
steps in there, it is as if with a single step he had come to a distant
place, infinitely far away from all noise and shouting and loudness,
from the terrors of existence, from the storms of life, from the scenes
of dreadful events or the languishing expectation of them. And wherever
you turn your gaze in there, everything will make you safe and calm.
The lofty walls of the venerable building, they stand so firm, they
guard so dependably the secure place of refuge, under whose mighty
vault you are free of every pressure. And the beauty of the
surroundings, their splendor, will make everything friendly to you, so
inviting; it will, as it were, endear the holy place to you. . . . .''

The reader understands that this attractive description is to be used
precisely in order to repel with it. ``How reassuring, how lulling,
alas, and how much danger in this security!''

Yes, but the speaker is nevertheless no fanatic who knows only how to
rail against something external, no iconoclast who reforms by smashing
things to pieces; nor is
% --- p. 13 ---
\opage{13}he an unloving man who can have no sympathy with those who, from
``noise and shouting,'' from ``the terrors of existence,'' and from
``the storms of life,'' seek a ``place of refuge'' behind the ``lofty
walls,'' an elevation of the soul under the ``mighty vault.'' It is a
religious thinker who in the discourse wishes to warn against the most
dangerous, most beguiling of all illusions, namely that of devotion,
when devotion is an aesthetic mood, a gentle illusion. ``How
reassuring, how lulling, alas, and how much danger in this security!
And therefore it is indeed truly so, that it is really only God in the
heavens who, in the actuality of life, can preach properly and to good
effect to human beings, for he has the circumstances, has the destinies,
has the strivings in his power. . . . . But in the house of God,
% TR: "i Guds Huus, i / i det pragtfulde Guds Huus": the print repeats "i" across the line break; rendered once.
in the splendid house of God, when the pastor preaches ---
for reassurance!'' . . . . . ``In earnest one wants to hear nothing of
the terrible; one wants to play at it, roughly as when in peacetime the
warriors, or rather the non-warriors, play at war; one wants,
artistically, to demand everything of the speaker, but oneself one
wants to sit, worldly and ungodly, quite secure in the house of God,
because one knows well enough that no speaker has the power, the power
that only Governance has, to seize a human being, to cast him into the
grip of circumstances, to let dispensations and testings and spiritual
trials preach to him in earnest for his awakening.''

Given the notions that attempts have been made from several quarters
to spread among the people, outsiders would now expect that the aim of
the discourse was, accordingly, to screw the severity of preaching up to
such a pitch that people would thereby necessarily be frightened away
from going to church. Anyone who has read the discourse knows how
loose, how superficial and false this notion is. The discourse is
indeed severe, but severe only insofar as the religious requirement is
in itself severe; there is no exaggeration in the discourse; here, as
everywhere in the religious, infinite severity is in turn essentially
mildness.
% --- p. 14 ---
\opage{14}``When faith'' --- thus the discourse ends --- ``is seen from its one
side, the heavenly side, one sees in it only the reflected radiance of
blessedness; but seen from its other side, the merely human side, one
sees sheer fear and trembling. But then that discourse is surely also
untrue which continually, and never otherwise than invitingly,
alluringly, winningly, will speak of the visit to the house of the
Lord, for seen from the other side it is terrifying. Therefore that
discourse is also untrue which ended at last by altogether frightening
people away from coming to the house of the Lord, for seen from the
other side it is blessed, one day in God's house better than a
thousand elsewhere''. . . . . ``For to the speaker it is said: use all
the ability granted you, willing to make every sacrifice and concession
in self-denial, use it to win people --- but woe to you if you win them
in such a way that you leave out the terror; therefore use all the
ability granted you, willing to make every sacrifice in self-denial, use
it to terrify people --- but woe to you if you do not, at bottom, after
all use it to win them for the truth.''

% TR: "Existensmeddelelse": "existence-communication" (term of art, first occurrence in this slice).
Since Christianity, according to Kierkegaard's conception, is not what
one simply calls a doctrine, but in the essential sense ``an
existence-communication,'' it was of particular importance to him to
make it as clear as possible how the speaker's own existence, without
the severity of the requirement giving the human a pretext to
evade the obligation and, into the bargain, blithely appeal to
``faith'' and to ``the pure doctrine,'' could and should relate itself
to the ``existence-communication.'' In this respect even the different
ways in which he had his writings published acquire a more than
accidental significance. His work ``\emph{For Self-Examination}'' he
published, e.g., under his own name, whereas ``\emph{Practice in
Christianity}'' is, in its first edition, ascribed to a pseudonym.
This is quite characteristic of what he calls his ``tactics.'' To the
language that ``Anti-Climacus'' uses in
% --- p. 15 ---
\opage{15}``Practice'' there corresponds so strenuous an existence that the
author, for that very reason, has placed himself behind the pseudonym
and announced himself only as editor: ``In this work, dating from the
year 1848, the requirement for being a Christian is forced up to a
highest pitch of ideality. Yet the requirement surely ought to be said,
presented, heard; Christianly, there ought to be no haggling over the
requirement, nor ought it to be suppressed --- instead of making
admission and confession concerning oneself.'' The work, on the other
hand, on which he has directly printed his name as author is meant, by
the spirit and tone in which the presentation is held from first to
last, to remind the reader of the relation of distance in which he
himself stands to the ideals.

% TR: "höiærværdig": the honorific of a bishop ("Right Reverend"); rendered "right reverend" to keep the allusion to Mynster.
``Let me,'' he says (For Self-Examination, p.~16), ``determine
precisely where, so to speak, I am. There is among us a right reverend
old man, the highest clergyman of this Church; what he, what his
preaching has willed, that same thing is what I will, only a tone
stronger, which lies in the difference of my personality, and which the
difference of the times demands''. . . . . ``Essentially I belong to
the average, and it is here that I have worked for unrest in the
direction of inwardization''.

``For Christianly there are two kinds of true unrest. The unrest in the
heroes of faith and the witnesses to the truth, which aims at reforming
an established order. So far out I have never ventured; that is not
something for me, and insofar as anyone in our time might seem to want
to venture so far out, I would not be disinclined to take up polemic
against him, to help make it manifest whether he is the one who has the
right''. . . . .

% TR: "umyndig Digter": lit. "a poet not of age" (legally under guardianship); rendered "without authority", Kierkegaard's own category.
``Instead of emptily and pompously passing myself off as a witness to
the truth and inducing others foolhardily to want to be the same, I am a
poet without authority, who moves others by means of the ideals.''

So much for the ``authorship'' in the first period. If the unbiased
observer now succeeds, by these
% --- p. 16 ---
% TR: "sikker Dom": "sure judgment", as quoted from the opponent (cf. the epigraph); to be harmonized with slice 1.
\opage{16}and similar features, in forming a tolerably ``sure judgment''
about S.\ Kierkegaard's fundamental thoughts and endeavors in this first
period of his public life, then a ``sure judgment'' about his ``mental
state'' in the last period will follow as of itself. As the premises
are, so too are the consequences, that is to say, when one has before
one a man who wills, whatever it may cost, to be consistent. To grant
the premises is therefore at once to acknowledge the consequences, that
is to say, if the granting rests on insight, not on a fleeting mood.

% TR: "Conseqvens" means both "consequence" and "consistency"; rendered by sense at each occurrence (Conseqvenserne = the consequences; conseqvent, Conseqvensen = consistent, consistency).
Now if it were always so entirely straightforward a matter that the
agreement of the consequences with their premises must at once strike
the eye of even a superficial observer, fuller explanations would here
be superfluous. But anyone who knows something of life's conflicts
knows within himself that the case occurs not so very rarely in which
what at bottom is pitiful inconsistency takes on the appearance of a
consistent stance, while conversely, what by appearance looks like a
break with all consistency is in its essence the only consistent thing.
The more inwardness the matter has, the more painfully will the
collisions be felt; if what is at issue is the problems of inwardness
themselves at their maximum, the knot will soon be drawn so tight that
the famous passage in Jacobi may come true in earnest: ``Ja, ich
bin der Atheist und Gottlose, der, dem Willen, der Nichts will, zuwider,
lügen will, wie Desdemona sterbend log. . . . .''

As the connecting link between the first and the last period of our
author's public life we have ``Practice in Christianity.'' To the new
printing of this work K.\ (Fædrel.\ 16 May 1855, no.~112) has sent a
guiding word. ``This book I have,'' he says, ``let appear in an
entirely unaltered printing, because I regard it as a historical
document. . . . . My earlier thought was: if the established order is
to be defended, this is the only way. . . . . In that way
% --- p. 17 ---
\opage{17}truth does after all come into the established order: it defends
itself by judging itself, it acknowledges the Christian requirement, it
makes confession concerning itself about its distance. . . . . This was
in my thinking the only way out for defending the established order
Christianly, and in order not to be guilty in any way whatever of
proceeding too hastily, I ventured to give the matter that turn, in
order to see what the old Bishop would do about it. If there was
strength in him, he had to do one of two things: either declare himself
decidedly for the book, venture to go along with it, let it be the
defense that averts what the book poetically contains, the charge
against the whole of official Christianity that it is an optical
illusion `not worth a pickled herring'; or else, as decidedly as
possible, throw himself against it, brand it as a blasphemous and
profane attempt, and declare official Christianity to be true
Christianity. He did neither; he did nothing; he merely wounded himself
on the book, and it became clear to me that he was powerless.''

% TR: "Skjöndt Beskyldningerne ere mange, er Skylden dog kun een": the Danish plays on Beskyldning (accusation) / Skyld (guilt) / for ... Skyld (for the sake of); only partly reproducible.
This is not the place to rebut all objections of every kind that have
been raised against S.\ Kierkegaard of late. Though the
accusations are many, the guilt is nevertheless only one: consistency!
For consistency's sake we shall here confine ourselves to two prominent
points: ``the attack on Bishop Mynster,'' and the attack on the
established order.

When Dr.\ Sören Kierkegaard, on Monday, 18 December 1854, with his
article\footnote{The article was written in February 1854.} in no.~295
of ``Fædrelandet,'' began with the question: \emph{Was Bishop Mynster a
``witness to the truth,'' one of ``the true witnesses to the
truth'' --- is this the truth?} there were certainly many who took
offense, many who cried inconsistency, many who thought of ``Herostratic
fame by hurling insults,'' etc., but, as it
% --- p. 18 ---
\opage{18}seemed, not very many who stopped to consider that the sensational
protest, in essentials, as far as the late Bishop is concerned,
contained nothing new, contained nothing other than what K.\ had already
said beforehand, if in a milder tone, to Bishop Mynster himself, to his
face, publicly and so plainly that the person concerned could not
misunderstand it, and that without either the Bishop himself or any of
his adherents at that time so much as making a show of having anything
to object against it.

The public and definite protest against Bishop Mynster's being assigned
a place among ``the true witnesses to the truth'' is to be read in
the work ``\emph{For Self-Examination},'' pp.~16 and 17. For after K.,
on p.~16, with the words already quoted: ``There is among us a right
reverend old man, the highest clergyman of this Church; what he, what
his preaching has willed, that same thing is what I will, only a tone
% TR: "min Persons Forskjellighed" here, "min Personligheds Forskjellighed" at the other two quotations (pp. 15, 18): rendered "my person" / "my personality" as printed.
stronger, which lies in the difference of my person'', has, in
expressions so telling that no one can mistake who ``the right reverend
old man, the highest clergyman of this Church'' is, compared his
activity with Bishop Mynster's, and thus directly placed himself at his
side, indeed, with the addition: ``a tone stronger, which lies in the
difference of my personality'', even --- for in such matters it is a
question neither of politeness nor of insult, but of truth --- if one
will, at his right hand; it is a matter of course that every word with
which K.\ characterizes his distance from ``the heroes of faith'' and
``the witnesses to the truth'' necessarily also comes to apply to the
Bishop. When, then, he says of himself that he ``essentially belongs to
the average,'' he says the same of Bishop Mynster; when he says of ``the
unrest in the heroes of faith and the witnesses to the truth, which aims
at reforming an established order'': ``So far out I have never
ventured, that is not something for me, and insofar as anyone in our
time might seem to want to venture so far
% --- p. 19 ---
\opage{19}out, I would not be disinclined to take up polemic against
him''. . . .; then anyone who has learned to read must involuntarily
understand what is said as if it read: ``So far out we, Bishop Mynster
and I, have never ventured, that is not something for us; and insofar as
anyone in our time might seem to venture so far out, we, Bishop Mynster
and I, would not be disinclined to take up polemic against him.'' In
short, every time Kierkegaard, with his ``admissions'' and
``confessions,'' bows before ``the heroes of faith'' and bends before
``the witnesses to the truth,'' it looks, in the reader's eyes, for all
the world as if he were holding Bishop Mynster by the hand, in order
thereby to induce this ``highest clergyman'' of the Church to do the
same.

Now if either the Bishop himself or any of those closest to him had felt
convinced that an encroachment had occurred here through this
comparison, why was this not said to K.\ at once and straight out? One
% TR: "Følgeblade til: Dette skal siges": "Følgeblade" (supplementary leaves) is Kierkegaard's heading for the dated sections (9 and 11 April 1855) that follow the main text of "Dette skal siges; saa være det da sagt" (1855), where the quoted passage stands; Nielsen's short title given in English.
could, after all, for as K.\ later said (Supplementary Leaves to: This Must Be
Said) ``in such transactions one uses words truthfully, like a
descriptive naturalist; compliments have no place here,'' have told him
without beating about the bush: ``By comparing yourself in this way with
`the highest clergyman of this Church,' it is not Bishop Mynster you
insult, but you violate the truth, in that by so unfitting a comparison
you confuse the concepts. You want to clarify existences; what, pray, is
your existence? You may well be the author with `the rich cultivation,'
`the deep religious foundation,' `the mind that is strong enough to look
life's highest problems in the eye,' but your existence? Yes, there you
are right, it `belongs to the average.' But Bishop Mynster's existence,
marked as it is with the signs of burdens, of crosses, of sufferings,
does not belong to the average; or can you perhaps, in selfish
blindness, not see that Bishop Mynster, besides what he is in the humane
direction, is also something other and higher than the men we otherwise
count among the foremost of the land, that he
% TR: The slice ends mid-sentence, inside the quotation begun at "Ved saaledes" (closed on p. 20). Full Danish sentence: "Men Biskop Mynsters Existens, saaledes mærket som den er, med Byrdernes, Korsenes, Lidelsernes Tegn, tilhörer ikke Gjennemsnittet; eller kan Du i selvsyg Forblindelse maaskee ikke see, at Biskop Mynster, foruden hvad han er i den humane Retning, tillige er et Andet og Höiere, end de Mænd, vi ellers regne blandt Landets Ypperste, at han just er et Sandhedsvidne, et af de rette Sandhedsvidner, et Led i den hellige Kjæde?“" Rendering of this part: "But Bishop Mynster's existence, marked as it is ..., does not belong to the average; or can you perhaps, in selfish blindness, not see that Bishop Mynster ... is also something other and higher than the men we otherwise count among the foremost of the land, that he".
% --- p. 20 ---
% TR: The slice opens mid-sentence. Full Danish sentence (p. 19--20): „eller kan Du i selvsyg Forblindelse maaskee ikke see, at Biskop Mynster, foruden hvad han er i den humane Retning, tillige er et Andet og Höiere, end de Mænd, vi ellers regne blandt Landets Ypperste, at han just er et Sandhedsvidne, et af de rette Sandhedsvidner, et Led i den hellige Kjæde?“ This slice renders from „just“: "is precisely a witness to the truth, one of the true witnesses to the truth, a link in the holy chain?''" (continuing "...that he").
\opage{20}is precisely a witness to the truth, one of the true witnesses to the truth,
a link in the holy chain?''

Had one been so openhearted from the first, then
K.\ would of course have known at once where he stood; now, on the other hand,
since silence was observed, he surely never dreamed that
it was perhaps only a diplomatic silence, but took,
as one sees, the silence straightforwardly for a tacit
admission.
% TR: „Indsigelsen“ = "the protest": Kierkegaard's article of 18 Dec. 1854 (Fædrelandet No. 295), which Nielsen calls an „Indsigelse“ on pp. 18 and 29 as well.
``Thus,'' he says in the protest, ``when
one lays the New Testament alongside, Bishop Mynster's
proclamation of Christianity was, especially for a witness to the truth,
a dubious proclamation of Christianity. But then there was, I
thought, this true thing in him, that he, as I hold myself
fully assured, was willing to confess before God and
himself that he was by no means, by no means a witness to the
truth---precisely this confession was in my thought the
true thing.''

Where, then, is the inconsistency? Had K.\ earlier
begun by complimenting Bishop Mynster as a ``witness
to the truth'' because at that time, for certain reasons, he liked
Bishop Mynster so well, but afterwards raised a ``Herostratic''
protest against calling Bishop Mynster a witness to the
truth because he now felt offended and could no longer
like Bishop Mynster: that would have been inconsistent.
Had he, moreover, to cover his inconsistency,
juggled with the word: witness to the truth, saying that it could of course, e.g., mean
both this and mean that, that he had admittedly
earlier happened inadvertently to use the expression:
witness to the truth, but had in no way meant
that one should thereby think of: the true witness to the truth,
whereas now, when he precisely wanted to point out what one should
think of by the true witness to the truth, Bishop
Mynster failed the test: then such conduct would not
only have been inconsistent, it would have been contemptible.
Now, on the other hand, since it is known before God and every man
that Kierkegaard himself, in the last as well as in the first period of
his life, holds fast to the same, the very same concept of
% --- p. 21 ---
\opage{21}``the witness to the truth,'' and, with respect to being a witness to the
truth, gives us in outline the same, the very same
picture of Bishop Mynster: in what, then, does the
inconsistency consist? What is it in this procedure of his
that is supposed to be apt to arouse suspicion against his
``mental state''?

How Kierkegaard wanted the assurance given in his cited
work (For Self-Examination p.\ 17) to be understood---that,
insofar as anyone in the present age might seem to want to venture ``so
far out'' (namely, like the ``witnesses to the
truth,'' to reform an established order), he was not
disinclined ``to take up polemic against him''---he has
shown in his article (Fædrel.\ 31 Jan.\ 1851 No.\ 26) against
Dr.\ Rudelbach. The article begins thus: ``The remark is
found'' (in Dr.\ Rudelbach's work on civil marriage,
1851) ``p.\ 70. In the text it says: Truly, the deepest
and highest interest of the Church in our days is. . . . precisely to be
emancipated from what can rightly be called habit- and
state-Christianity. To this, then, note 121. This is the
same thing that one of our most distinguished writers of
recent days, Sören Kierkegaard, seeks to impress,
to imprint, and, as Luther says, to drive into all those who
will hear.''

This article is, in accordance with Kierkegaard's
principle, composed in a strictly conservative spirit. ``Nothing,''
it says, ``makes me so uneasy as everything
that even merely smacks of this unhappy confusion
of politics and Christianity, a confusion that can so easily
bring a new kind of church reformation up and into fashion,
the backward reformation which, reforming, puts a
new worse in place of an old better, while it
is nonetheless supposed to be certain and true that it is a reformation,
on which occasion the whole town is illuminated. Christianity
is inwardness, inwardization. If at a given
time the forms under which one is to live are not the most
perfect: well then, in God's name; if one can get them better:
% --- p. 22 ---
\opage{22}good. But essentially: Christianity is inwardness.'' . . . .
``Should this belong to true Christianity, this faith
in the saving power of politically attained free institutions,
then I am no Christian, indeed, worse still, I am
a very child of the devil; for, candidly speaking, I am even
suspicious of these politically attained free institutions,
especially of their saving and regenerating power. That is
% TR: „christendum“ is a printer's error, the word-space dropped from „christen dum“ (as the transcriber notes); rendered by that sense, "so stupid a Christian".
my Christianity, or so stupid a Christian am I, who
moreover have never concerned myself with `Church' and `State'
---that is much too big for me . . . . Nor have I
ever fought for the emancipation of `the Church,' any more than
for that of the Greenland trade, of women, or any
other emancipation whatsoever. As a single individual, with
the aim `the single individual,' aiming at inwardization in Christianity
in `the single individual,' I have consistently, with the weapons of the spirit,
solely and only with the weapons of the spirit, fought to call
% TR: The print has only one closing mark after „Sandsebedrag“ for both the inner quotation and the quotation opened at „Skulde“; the English closes both.
attention to, and against letting oneself be deceived by, `sense-deception'\,''
. . . . ``It is not externally that Christianity
is to be helped, by institutions and constitutions, and then
least of all when these are not, in the old Christian way,
to be won through suffering by martyrdoms, but, in the political way,
socially-amicably gained by balloting or in
the number lottery---to be helped in that way is on the contrary
its ruin.''

If, however, anyone thinks he finds in these utterances of a
``conservative'' disposition a proof that
Kierkegaard had thereby so directly signed himself over
to the established order that he could never, without expressly breaking
with his own principles, permit himself
to attack the established order, such a one will, by
looking a little more closely, easily convince himself that he has
been remarkably mistaken.

``Let me, however,'' it says, as in a postscript to
the article, ``add the following, so that what I say may not
be misunderstood, as if it were my opinion that
Christianity consisted solely and only in putting up with everything
% --- p. 23 ---
\opage{23}concerning the external forms'' . . . . ``In the Acts of the
Apostles we read the words: one ought to obey God rather than
men. So there are cases where an established order
can be of such a nature that the Christian ought not to
put up with it, ought not to say that Christianity is precisely
this indifference to the external.''

``But let us now see how the Apostles did not conduct
themselves---for how these venerable figures did conduct
themselves, surely everyone knows.'' ``The Apostles did not go about
like that chatting with one another and saying: `It is
intolerable that the Sanhedrin imposes punishment on the proclamation
of the Word; it is coercion of conscience. Yet what are
we to do? Should we not try to get a few of us together, and
then submit an address to the Sanhedrin, or try to get
onto a synod; it would not be impossible that we then,
by holding together with what are otherwise our enemies,
could by balloting get the majority, and we could then
get freedom of conscience to proclaim the Word.' God in
heaven! Venerable figures, forgive me that I have had
to speak thus; it was necessary.''

``How, on the other hand, did they conduct themselves---for there are
surely a good many who have forgotten it.'' . . . . ``The Apostle
is essentially a solitary man.'' . . . . ``Let us suppose
that someone meets him who says: Do you know that the Sanhedrin
has made scourging the penalty for proclaiming the Word. The Apostle
answers: Well, if the Sanhedrin has done so, then I shall accordingly
be scourged. Tomorrow the Sanhedrin imposes the death penalty.
The Apostle answers: Well, if the Sanhedrin has done that, then I shall
accordingly be executed. He lets the established order
stand; not a word, not a syllable, not a letter
in the direction of change in the external, not the most fleeting
thought in his head, not a blink of the eye, not the
movement of a feature in this direction.
% TR: The print opens a quotation at „Lad os antage“ and inner ones at „Nei,“ and „lad det“, but closes only once after „Amen.“; the English closes the inner and the outer quotation there.
`No,' says the Apostle,
`just let it stand unshaken, this established order, for
this too stands unshaken, with God's help, that today
I am scourged, tomorrow executed, or, what
% --- p. 24 ---
\opage{24}is the same thing, today I proclaim the Word, and tomorrow,
Amen.'\,''

% TR: p. 24 "Ordenligviis" (no t) is so printed (checked at the image); p. 27 quotes the same sentence as "Ordentligviis". Rendered "Ordinarily".
% TR: „Selvbekymring“ = "self-concern"; „eenlig Mand“ = "solitary man" throughout.
``And here a Christian memento. Ordinarily,
Christianity within Christendom is precisely indifference to
the external in self-concern. But if someone so
collides with the established order that it could occur to him
that it is a question of conscience---great God, a
question of conscience!---and he ventures to say this:
then he has to be a solitary man, to fight in suffering, to
choose martyrdom. Conscience and a matter of conscience can
---this is eternally certain and lies in the matter itself, for
conscience is no numerical determination---only be
represented by a solitary man and in character, through
action.'' . . . .

When this article---precisely an article after Bishop Mynster's
own heart---saw the light, there was hardly anyone among
those who then held to the established order who
found anything offensive in it, let alone anything so offensive
that it could in the remotest way arouse suspicion against
the author's ``mental state.'' The principles stated in the article
are indeed also really---as everyone
who will use his understanding must admit---of such a nature that
both the State and Christianity, that is to say, the Christianity of the New
Testament, can, each according to its own interest,
acknowledge them. What, in religious respects,
causes the State its many vexations, clashes and disagreeable
troubles is by no means ``the solitary man who chooses
to fight in suffering,'' but much rather the politicizing
church system, the religious worldliness that nevertheless wants to rise above
the worldly,
% TR: „Selskabsdrift“ = "gregarious instinct" (the social or herd drive), here and on p. 26.
that gregarious instinct of consciences which
founds its freedom on ``balloting.'' ``Everything that
is party and wants to act as party, perhaps even intriguingly
---when it wants to appeal to `conscience' against an established order,
it incurs the guilt of an untruth,'' is
a principle which, remarkably enough, can be embraced in equal measure, even if
in opposite spirit, by statesmen and apostles.

% --- p. 25 ---
\opage{25}If we now ask whether Kierkegaard's, as it seems,
in the highest degree agitating, violent and revolutionary
public conduct in the last periods of his life can, without circumlocution,
excusing turns of phrase or subtle artifices,
be brought into agreement, into thorough, true and
complete agreement, with the train of thought in the just-mentioned
peaceable article, the answer is not far off.
Only we must not, as regards the last scenes, forget
that the bow is drawn, not forget what
a high-flown ``testimony'' it was that gave occasion for
the bow's being drawn thus, not forget that the knot of
inwardness has been pulled tight, that it is now earnest, really
earnest with ``the famous passage'' in Jacobi: Ja, ich bin
der Atheist und Gottlose!

% TR: The print never closes this quotation (no mark before the parenthesis); the English supplies the closing marks after the letterspaced span.
% TR: „Anden Pintsedag“ = Whit Monday.
Example: \emph{``What is to be done---let it be done
now through me or through another.}'' (Whit Monday
1854.\footnote{Here may be repeated what K.\ himself expressly stresses in an
earlier article (Ved Biskop Mynsters Död): ``Note
\emph{the year and the date}.''} S.\ Kierkegaard. Fædrelandet, 22 March 1855.
No.\ 69.)

``First and foremost, on the largest possible
scale, an end must be made of the whole official---well-meant---untruth
which---well-meant---conjures up and
maintains the illusion that it is Christianity that is proclaimed,
the Christianity of the New Testament. In this respect
the point is not to spare.'' . . . .

``When that has been done, the matter is to be turned thus:
is not the true state of affairs really this, that
from generation to generation things have gone so far backward for
us human beings, . . . . that we, instead of the impudent
drivel about Christianity being perfectible, about our
going forward, indeed, about Christianity perhaps no longer
even satisfying---that we, wretched, pitiful dozen- and
score-men, in the end really
% --- p. 26 ---
\opage{26}cannot bear this divine thing which is the Christianity of the New
Testament; and that we must to that extent content ourselves
with the kind of religiousness that is now the official one, after
it has, be it noted, been made recognizable that it is by no means
the Christianity of the New Testament. . . . Thus must
% TR: „Öienforblindelse“ = "blinding of the eyes", i.e. self-deception.
the matter be turned; away, away, away with all blinding of the eyes,
forth with the truth, out with it: we are not fit
to be Christians in the New Testament's sense,---and
yet we need to dare to hope for an eternal salvation.'' . . .
% TR: „Styrelsen“ = "Governance" (Providence), capitalized as in Kierkegaard's usage.
``When the matter has then been turned thus, it will become evident
whether there is anything true in this, whether it has Governance's
consent,---if not, then everything must burst, so that in
this horror individuals might again come to be who can
bear the Christianity of the New Testament. But an end,
an end must be made of the official---well-meant---untruth.''

But is it not, after all, outrageous? ``Ja, ich bin der
Atheist und Gottlose, der dem Willen, der Nichts will, zuwider,
lügen will.'' Had K., instead of engaging in
such ``radical'' demands and ``pessimistic'' views,
been able to bring himself to strike out on a smoother road;
had he, tired of his position as ``solitary man,'' at
last followed the Christian gregarious instinct, or, to choose
a better expression, been so ``loving'' as to be like others
and act like others: things would have gone quite
differently for him. Had he, instead of inveighing against ``official
Christianity'' and causing scandal, shown the
prudence to come forward, e.g., with a proposal for the ordering
of the expected church constitution: he would surely at some
point have succeeded ``in getting onto a
synod''; and once he had got onto the synod, and
had been given the floor at the synod, perhaps even the most
votes in the ``balloting,'' then surely no cautious
man would have ventured to say aloud of S.\ Kierkegaard
that he was now acting inconsistently, or have permitted himself to raise
any objection against his ``mental state.''

% --- p. 27 ---
\opage{27}The smooth road was inviting enough; why, then, did
he not choose the smooth road?

To dazzle the thoughtless is a poor art; but
anyone who will see must be able to see that K.\ could not, without
breaking with his presuppositions, engage in
reforms by the peaceful road. For him there was only a choice
between one of two things: either to remain standing, as hitherto, on
the proposition: ``Ordinarily, Christianity within
`Christendom' is precisely indifference to the external in
self-concern,'' or, when the course of things took this
``ordinarily'' away from him, then to attack the established order
as resolutely, as unconditionally, as decisively
as possible.

``As carefully as it has hitherto been concealed
what my task might become, as cautiously
as I have been preserved in the most impenetrable
unrecognizability: just as decisively shall I also now, when
% TR: „da Öieblikket er kommen“: "now that the moment has come", with a play on the title of the periodical Øieblikket (The Moment); likewise „Öieblikkene“ below (plural: the issues of The Moment), rendered ``the Moments''.
the moment has come, make it recognizable. The question
what Christianity is, and with it in turn the question
of the State Church, whatever one wants to call it, the combination
or the holding together of Church and State, shall be
brought to the utmost decision. It cannot and shall not
go on as it went on year after year.'' . . .

But is it not, then, open revolt? When K., in
the work ``For Self-Examination,'' p.\ 16, asks: ``Am I then
perhaps proclaiming tumult, the overturning of everything, disorder?'' and
then answers: ``Truly, no''; could one not
with regard to the ``Moments'' turn his words as a
weapon against him, and say: ``Truly, yes. Now you are proclaiming
tumult, the overturning of everything, disorder!''? This matter
is not complicated. The Kierkegaardian language is, without
circumlocution, revolutionary. ``For is it not true, you common
man, that you can understand very well---I mean that
precisely you can understand it far more easily and better than
demoralized pastors and a corrupted gentility---is it not
true, this you can understand excellently, that it is one thing
to be persecuted, mistreated, scourged, crucified, beheaded
% --- p. 28 ---
% TR: „o. D.“ (og Deslige) = "and the like", here and („o. dsl.“) on p. 29.
\opage{28}and the like, another thing, comfortably settled, with a
family, steadily advancing, to live by depicting that another
was scourged, etc.'' This language is revolutionary;
no agitator will easily outbid it.

Nevertheless, Kierkegaard's public conduct is the complete opposite of the
revolutionary; we have incontrovertible proofs
of this. ``If someone so collides with the
established order that it could occur to him that it is a
% TR: „støre Gud“ is a printer's error (ø sort) for „store Gud“, as the transcriber notes (cf. p. 24); rendered "great God".
question of conscience---great God, a question
of conscience!---and he ventures to say this, then he has
to be a solitary man, to fight in suffering, to choose
martyrdom.'' These words Kierkegaard conscientiously strove,
during his last public conduct, to put into character;
we may well stress this fact without therefore
needing to speak in high-flown terms and compliment
the deceased as a martyr, or as a ``witness to the
truth.'' K.\ was and remained throughout the whole struggle a ``solitary
man'': there was no one who helped him, no one who could,
in a finite sense, come to his help. Even if there
had been---which, however, there probably was not---someone or other
among us who, with the same sureness about the categories,
had also had the resoluteness and courage to venture
equally far out, it would still not have come to a finite
alliance between the two: each of the combatants,
even if they had fought in the best understanding in the same
spirit, for the same cause, would still have had to fight on his own,
as a solitary man. But that such a procedure
is the exact opposite of the revolutionary one, surely
everyone sees. If one overlooks this contrast, how
will one then explain to oneself that Kierkegaard's revolutionary
language could be heard by ``the common man''
without anyone in the crowd stirring so much as a finger
to attack the external? Or did Kierkegaard's
words perhaps not know how to find the way to the passions?
Does his language perhaps lack the inflaming quality that
otherwise so easily produces a ferment? Or could ``the
% --- p. 29 ---
\opage{29}common man,'' in his simplicity, nevertheless not
quite understand him when he said ``that it is one thing to be persecuted,
mistreated, scourged, crucified, beheaded, and the like,
another thing, comfortably settled, with a family, steadily
advancing, to live by depicting that another was
scourged?'' Or was it perhaps the apostolic authority of the pastors' church
that impressed the common man?

For the fact that the Kierkegaardian protest, the strongest that in
recent times has been made against an established church system, a
protest which, while it sounded the whole country over with an echo
in thousands upon thousands of souls, nevertheless first and last sounded
in such a way that every man understood that what it was about was not an
overturning in the external but a change in the internal,
for this we have to thank not only the people's
level-headedness and sound sense---sound sense
could, alas, in such things easily be beguiled by an agitator---but
essentially and categorically Kierkegaard's position as
``solitary man.''

What has otherwise been achieved here by the Kierkegaardian
protest cannot be made clear in this context.
% TR: The print has a stray full point before the comma after „Forudsætninger“ (printer's error, noted by the transcriber); not reproduced.
The task was to point out in some features the consequences
of his presuppositions, the inner agreement
of his word and action from first to last. Those
who lament that they cannot have ``any sure judgment''
about this man's ``mental state in the last period
of his life'' ought seriously to consider whether the reason for this
might not, after all, lie with themselves, namely that
% TR: „komme“ is a printer's error for „komne“ (noted by the transcriber); rendered "have come".
in their researches they have not yet come so far
that they can have any sure judgment about ``this man's''
authorship in the first period.

% End of the article: a short centered rule; no signature or dateline.
\begin{center}\rule{0.25\linewidth}{0.4pt}\end{center}

\end{document}
